\documentclass[11pt]{article}
\usepackage[margin=1in]{geometry}
\usepackage{amsmath,amssymb}
\usepackage{hyperref}
\usepackage{listings}
\usepackage{xcolor}
\usepackage{enumitem}

\lstset{
  basicstyle=\ttfamily\small,
  breaklines=true,
  frame=single,
  columns=fullflexible,
  keepspaces=true,
  showstringspaces=false
}

\title{AuditPledge: Provenance and Priced Risk for the MCP Tools\\ a Compound-Facing Agent Calls}
\author{Dipankar Sarkar \\ \small doom2quake}
\date{August 2026}

\begin{document}
\maketitle

\begin{abstract}
An AI agent that manages a Compound position is only as trustworthy as the weakest tool in
its toolbelt. Agents increasingly acquire those tools over the Model Context Protocol (MCP),
and one unsigned, remote-code-capable, over-scoped tool can read a signing key, rewrite a
transaction before it is broadcast, or move collateral. The operator wiring such an agent has
no standard way to answer which of the tools it can call are dangerous, how dangerous, and in
money. We present AuditPledge, a supply-chain auditor that opens a real MCP session against an
agent's tool servers, verifies Ed25519 provenance for every tool against a committed public
keyring, scores the risk weighted for what a Compound position exposes (state reads,
transaction building, collateral movement), and prices the residual risk into an expected loss
and a suggested premium. A counterfactual pass reruns the identical scorer over a mutated
manifest, so the dollar saving from a fix is measured, not asserted. Everything ambiguous fails
closed: unknown servers, identity mismatches, forged or missing signatures never read as clean.
The audit, the cryptography, the scoring and pricing, and the UI numbers are real; the demo
tool catalogue, the offline advisory snapshot, and the demo publisher keys are synthetic and
labelled as such. The deterministic core is verified by 75 Python tests, needs no network, key,
or cloud credentials, and all Web3 interaction is testnet and seeded only. This is a candidate
entry to the Compound Grants Program, Security Tooling domain (Section~\ref{sec:limits}).
\end{abstract}

\section{Problem}
DeFi is being wired up to AI agents: position managers, liquidation bots, rebalancers, and
treasury assistants increasingly run as agents that call external tools over
MCP~\cite{mcp}. MCP has become the default way an agent acquires tools, ahead of any security
convention for it. A tool manifest is attacker-controlled data. A tool that merely
\emph{claims} a publisher and \emph{claims} a signature proves nothing, yet an operator today
reads a list of tool names, trusts the publisher on faith, and finds out something was wrong
only after funds move.

Software supply-chain attacks --- typosquatted packages, compromised maintainers, lookalike
publishers --- are already a dominant way real systems get breached~\cite{osv}. The MCP tool
layer imports every one of those failure modes and adds a live wire to a wallet. The failure
sharpens under a Compound-facing agent, because the tools that matter are exactly the ones
that sit in front of the money: readers that pull Comet~\cite{comet} market and account state,
builders that assemble and sign COMP-related transactions, and movers that supply or withdraw
collateral. The concrete gap is the absence of three things at once: a provenance check on the
tool supply chain, a severity ranking weighted for Compound exposure, and a dollar figure a
risk owner or an insurer could act on.

\section{Why Compound}
Compound is a mature lending protocol with real value at stake and an active, security-minded
governance process~\cite{compgov} whose Grants Program names Security Tooling as a funded
domain. A security tool is not a one-shot deliverable; it is an artifact the ecosystem keeps
using, which is the shape milestone funding supports. AuditPledge is not a chain-agnostic tool
with a Compound label: its risk model weights capability against what a Compound position
actually exposes, and its reference keyring is seeded from publishers the Compound community can
vet. A provenance-and-pricing standard proven against Compound-facing agents is a public good
the whole ecosystem inherits, and the moment to set the norm --- that a tool a Compound agent
calls carries verifiable provenance --- is before an incident sets it the hard way.

\section{Design}
AuditPledge is a pipeline of pure functions over an MCP tool manifest and a committed keyring,
so a reviewer can trace every number. The stages are: inventory, provenance verification, risk
scoring, underwriter pricing, and counterfactual repricing.

\subsection{The evidence path: a real MCP session}
With \texttt{--source mcp}, AuditPledge opens an MCP stdio session, completes the
\texttt{initialize} handshake, checks the server identity it was handed against the server it
asked for, and reads \texttt{tools/list} over JSON-RPC. Unknown servers, identity mismatches,
and unreachable servers all \textbf{fail closed}: the audit stops rather than reporting a clean
manifest. A bundled synthetic Compound tool catalogue is read the same way over the same
session, and is labelled \texttt{synthetic: true} everywhere it surfaces.

\subsection{Provenance: Ed25519 over canonical bytes}
Each tool carries a detached Ed25519~\cite{ed25519} signature over the canonical bytes of a
fixed field set:
\[
\text{signed} = (\text{name},\ \text{version},\ \text{publisher},\ \text{sorted scopes},\ \text{sorted ops}).
\]
The verifier recovers nothing from self-assertion. It normalises the publisher and matches it
\emph{exactly} against the keyring --- \texttt{attacker.compound-finance} is not
\texttt{compound-finance}, because suffix matching is never used --- and then verifies the
signature over the canonical bytes with that publisher's committed public key. Only a signature
that verifies counts as signed; a present-but-invalid signature is recorded as \emph{forged}.
Tampering with any covered field (adding an op, widening a scope, bumping a version) invalidates
the signature. Every ambiguous path (no signature, unknown publisher, non-hex signature, crypto
library unavailable) returns \texttt{verified = false}, so the auditor treats the tool as
unsigned rather than guessing.

\subsection{Risk: weighted for Compound exposure}
Each tool's score is a sum of named weights over its declared capabilities and provenance. RCE-
class ops (\texttt{exec}, \texttt{run\_command}, \texttt{write\_file}, and critically
\texttt{sign\_tx}, since on a Compound-facing agent a compromised signer moves real value)
carry the heaviest weight; network ops, unsigned provenance, an untrusted publisher, and
over-broad scopes (four or more, or a wildcard) each add their own. A distinct \textbf{Compound
position exposure} weight is added when a tool can build or sign a transaction or move
collateral, because that is what sits directly in front of the money in a Comet position.
Advisory hits from a source that must name its ID, origin, and retrieval time add further.
Score bands map to \texttt{low/medium/high/critical}.

\subsection{Pricing and counterfactuals}
An underwriter pass converts the assessment into money. Per tool, a compromise probability
$p$ is read from its severity band, floored at a higher value when a tool is both unsigned and
RCE-capable (the dominant real-world failure), and the expected loss is $p \times$ a
configurable breach-loss assumption. Summed and multiplied by a risk load, that is a suggested
annual premium. The value of a fix is then \emph{measured}, not asserted: a counterfactual pass
reruns the \emph{identical} \texttt{assess} $\rightarrow$ \texttt{price} pair over a mutated
manifest --- the worst tool dropped, or re-published by a trusted publisher with its RCE/signing
ops and wildcard scopes removed and \emph{genuinely re-signed} with that publisher's key --- and
reports the premium delta. Because the mutated manifest passes the same Ed25519 verification as
any other, the repricing is a real rerun of the scorer, not a hand-edited number.

\subsection{The agent graph}
An optional layer turns the audit into a ranked remediation plan. It is an
ADK~\cite{adk} \texttt{SequentialAgent} that runs an Auditor sub-agent, then a Risk Lead
sub-agent, in that fixed order; a run that did not write both output keys into session state is
discarded rather than passed off as a plan. The order is enforced by the workflow agent, not
requested in a prompt. The deterministic audit is the whole security and pricing claim and runs
without any cloud credentials; the LLM layer only adds narrative on top of numbers the code
already produced.

\section{Implementation and evaluation}
Provenance verification uses a standard Ed25519 implementation~\cite{ed25519}; everything else
--- inventory, scoring, pricing, counterfactuals --- is dependency-light Python that is a pure
function of its inputs. The static UI (\texttt{auditpledge/ui/index.html}) carries no
hand-written numbers: every figure is captured from one real audit run and pasted between
snapshot markers, and a snapshot test re-runs the audit and fails if any value drifts.

We evaluate by test with pytest:

\begin{itemize}[leftmargin=1.4em]
\item \textbf{75 Python tests} (\texttt{PYTHONPATH=. pytest -q}, Python 3.10+, the audit store
forced in-memory). They pin each defence: the Ed25519 verification and its tamper cases,
exact-publisher matching, the unsigned-RCE pricing floor, the counterfactual repricing, the
advisory-ID validation (a placeholder such as \texttt{GHSA-xxxx-xxxx-xxxx} is refused rather
than printed as a known vuln), the fail-closed MCP paths against a real stdio subprocess
(handshake, \texttt{tools/list}, identity mismatch, missing evidence, unreachable server), the
agent-graph ordering, and the UI snapshot. No network, key, or cloud credential is needed.
\end{itemize}

The sharpest case is the hostile manifest \texttt{evil-lookalike}: a tool that claims the
\texttt{compound-finance} publisher with a bare-string ``signature'' and a lookalike
\texttt{attacker.compound-finance} entry. The auditor recovers no trust from either: the first
is recorded as forged, the second fails the exact-publisher match, and both are scored as
unsigned. That pins the property that self-asserted provenance is not trust.

\section{Limitations}\label{sec:limits}
This is milestone 1, and the honest scope is narrow (the full map is in
\texttt{docs/HONESTY.md}).
\begin{itemize}[leftmargin=1.4em]
\item \textbf{No mainnet, no users, no revenue.} All Web3 interaction is testnet and seeded
only; nothing is on Compound mainnet.
\item \textbf{No third-party MCP server has been audited by us.} AuditPledge demonstrates both
ends against its own synthetic catalogue. A real third-party tool does not yet publish
provenance evidence, which is exactly what the milestone roadmap delivers; against such a tool
the auditor fails closed rather than inventing a verdict.
\item \textbf{The pricing inputs are illustrative priors, not an actuarial model.} The
per-severity probabilities, the unsigned-RCE floor, the breach-loss figure, and the risk load
are defensible, runtime-configurable starting points. What is demonstrated is the mechanism: a
manifest becomes a reproducible number and a fix changes it by a measured amount.
\item \textbf{No audit, partnership, or endorsement} from Compound or anyone else. This is a
grant application.
\end{itemize}

\section{Conclusion}
AuditPledge turns ``which of my Compound agent's tools are dangerous, how dangerous, and in
money'' into a reproducible, provenance-checked, priced answer that fails closed on every
ambiguous path. The counterfactual makes the value of a fix a measured saving rather than a
claim, and the UI cannot show a number the auditor did not produce. The durable contribution is
a provenance-and-pricing discipline for the MCP tool supply chain that Compound-facing agents
depend on, established before an incident forces it.

\begin{thebibliography}{9}
\bibitem{mcp} Anthropic. \emph{Model Context Protocol Specification}. \url{https://modelcontextprotocol.io/}, 2024--2026.
\bibitem{comet} Compound Labs. \emph{Compound III (Comet) Documentation}. \url{https://docs.compound.finance/}, 2023--2026.
\bibitem{compgov} Compound Governance Forum. \url{https://www.comp.xyz/}, 2026.
\bibitem{ed25519} D.~J. Bernstein et al. \emph{High-speed high-security signatures (Ed25519)}. Journal of Cryptographic Engineering, 2012.
\bibitem{osv} Open Source Vulnerabilities (OSV). \emph{OSV.dev distributed vulnerability database}. \url{https://osv.dev/}, 2021--2026.
\bibitem{adk} Google. \emph{Agent Development Kit: Workflow Agents (SequentialAgent)}. \url{https://google.github.io/adk-docs/}, 2025--2026.
\end{thebibliography}

\end{document}
